\chapter*{\runtitulo}

\noindent
El problema de \emph{Community Membership Hiding} (CMH), u \emph{ocultamiento de
membrres\'ia de comunidades}, consiste en, dado un grafo, un algoritmo de detecci\'on
de comunidades y un nodo objetivo $u$, encontrar un conjunto m\'inimo de reconexiones
de aristas adyacentes a $u$ de modo que el algoritmo deje de asignarlo a su comunidad
original, mientras se preserva la integridad estructural del grafo.
Este problema es relevante en el contexto de la privacidad en redes, donde los
algoritmos de detecci\'on pueden revelar informaci\'on sensible sobre los usuarios---como
afiliaciones pol\'iticas o ideol\'ogicas---a partir de su posici\'on estructural, sin
que estos lo hayan divulgado expl\'icitamente.

Trabajos recientes formularon CMH como un Proceso de Decisi\'on de Markov (MDP) resuelto
mediante Aprendizaje Profundo por Refuerzo (DRL), empleando una arquitectura
Actor-Cr\'itico (A2C). En este trabajo identificamos una limitaci\'on arquitect\'onica
cr\'itica de dicho m\'etodo: el agente eval\'ua el grafo de forma global sin condicionar
expl\'icitamente en el nodo objetivo, lo que le impide aprender una pol\'itica de
ocultamiento espec\'ifica.

Proponemos una arquitectura DRL novedosa que incorpora expl\'icitamente la identidad del
nodo objetivo como caracter\'istica de entrada, utiliza \emph{Graph Attention Networks v2}
(GATv2) como codificador de grafos para aprovechar un mecanismo de atenci\'on din\'amico
y m\'as expresivo, y combina se\~nales est\'aticas y din\'amicas que reflejan los cambios
topol\'ogicos en tiempo real durante el proceso de reconexci\'on.
Extendemos el sistema de un \'unico agente a una \emph{Arquitectura Dual} que especializa
agentes separados para la adici\'on y la eliminaci\'on de aristas, e introducimos una
heur\'istica de \emph{Subgrafado Guiado por Pol\'itica} que hace el m\'etodo escalable
en grafos de gran tama\~no.

La evaluaci\'on experimental sobre los benchmarks est\'andar del \'area demuestra que
nuestra arquitectura supera consistentemente al agente DRL original y a los m\'etodos
basados en gradientes, logrando altas tasas de \'exito con una perturbaci\'on estructural
m\'inima.

\bigskip

\noindent\textbf{Palabras clave:} detecci\'on de comunidades, ocultamiento de
membrres\'ia de comunidades, aprendizaje por refuerzo profundo, redes de atenci\'on de
grafos, privacidad en redes, GATv2, Actor-Cr\'itico.
